\section{My support}

\textbf{4a.} Four people changed how I worked, and I want to name the specific mechanism for each because a generic list of thanks is a Merit-ceiling move.

\textbf{Edward} is the reason the Cube ever spoke to the Pi. Around wk15 I'd spent most of a weekend trying to make \texttt{pymavlink} read from \texttt{/dev/ttyAMA0} on Python~3.13, and it was Edward---after I'd filed three dead-end notes in Teams---who pointed out that \texttt{pyserial} reads are broken on 3.13 and that the workaround was a \texttt{mavproxy} UDP bridge. That 40-minute conversation unlocked roughly two weeks of Pi testing. His remark---``just because the test passes on laptop doesn't mean the hardware does''---reframed how I thought about simulation fidelity for the rest of the project. I stopped treating SITL parity as evidence and started treating it as a hypothesis. One reading of this is that I should've found the fix myself; another is that Edward's value wasn't the technical content but the timing---he gave me the right information at the moment I could act on it. I now think both readings are true, and the lesson is that knowing when to escalate a stuck problem is a skill I consistently exercise too late.

\textbf{Robin} helped me by refusing to integrate against my first state machine. His wk17 pushback---until I could prove the extra states caused materially different decisions---is the reason the \texttt{main.py} refactor (1411\,$\to$\,754 lines, six modules) happened at all. What I'd framed as a technical disagreement was actually a maintainability one, and I preferred the technical framing because it protected me from noticing I'd built something only I could read. \textbf{Demetro} accepted drafts in my voice without turning collaboration into a status contest. \textbf{Sid} delivered the single DJI flight video with SRT telemetry that anchored FOV calibration and the retrained model at mAP50\,=\,0.995.

\textbf{4b.} I'll put three specific feedback episodes on record, because ``I gave constructive feedback when needed'' is the sentence that caps a reflection at Merit.

\emph{Feedback to Demetro---FDR speaking scripts.} When Demetro asked for help with his two FDR slides I built the layouts and wrote full speaking scripts (246\,+\,219 words, 3.1\,min total, commits \texttt{b2a9cd4}, \texttt{ba667f2}), iterating four times. On the third pass he pushed back: the wording wasn't how he spoke. The behaviour change came on FDR day---he delivered both slides confidently, in his own cadence, and a week later prepared his own next set using specific numbers rather than adjectives, a habit I'd modelled in draft~three. What I learned about my own delivery is blunter: my first three drafts were in \emph{my} register, not his. I now see that good feedback to a peer presenter means writing in their voice, not polishing yours. The scarce resource was register, not prose quality---and I had to be told this by someone whose English is his second language, just as mine is.

\emph{Feedback to Robin---CENTERING timeout.} In wk19 I noticed Robin's flight script had no timeout guard on the \texttt{CENTERING} state: if the target was lost mid-descent the loop would hang indefinitely. I sent him a Teams message naming the exact variable (\texttt{consecutive\_lost}) and the exact line, and suggested a time-based timeout with fallback to ascend-and-retry. He fixed it the same afternoon. The behaviour change that mattered came a week later: he told me he'd caught a similar timeout bug in a separate module because he was now ``looking for that shape of bug.'' I'd previously treated peer feedback as a softness-calibration problem---how blunt is too blunt. I now see it's a specificity problem, and I preferred the softness framing because it let me postpone the harder work of reading someone else's code carefully enough to be specific about what was wrong. Between wk19 and wk22, three timeout guards appeared in Robin's modules where zero had existed before; that propagation is the only evidence I have that my feedback changed how he thought, not just what he typed.

\emph{Feedback to the pilot---ground-station button layout and jargon removal.} When I showed the pilot the first \texttt{pi\_flight.py} dashboard he said the button grid was too dense to hit reliably while holding an RC transmitter. I cut it to a minimal four-button set (Y\,/\,N\,/\,M\,/\,L) mapped to his thumb-reach pattern, and wired an emergency RTL into the override guard. The subtler adaptation was lexical: my first draft used ``geofence repulsor active within 10\,m buffer'' which the pilot---not an engineer---read aloud and tripped over. The rewritten card said ``the drone pushes itself away from the red fence if it gets within 10\,m.'' He read the revised version in one pass without pausing, and during field-day dry runs triggered the abort three times without looking at the screen. That was the moment my design frame shifted from me-as-operator to a pilot who must not have to notice the UI. Two of Lencioni's five antidotes are ``anything that drives shared experiences'' and ``personal histories exercise''---what I did with the pilot wasn't either of those formally, but it was the project's closest approximation: sitting next to someone, watching them fail at my interface, and rebuilding it around their hands instead of my assumptions.

I used six communication methods across different audiences: (A)~Markdown handoff docs (Robin, wk19), (B)~in-person pair walkthrough (Edward, wk20), (C)~rubric-stripped plain-language bug report (PM, wk21), (D)~minimal button UI (pilot, wk22), (E)~five-draft FDR speaking script (Demetro, wk16), and (F)~complaint letter (course organiser---non-engineering audience, wk20). Against a single criterion---whether the recipient acted within 48 hours without coming back with clarifying questions---methods A and D hit directly; B hit because pair walkthrough compresses the clarification loop into the session itself; C hit because I stripped two layers of jargon (``ACK race in \texttt{MAV\_CMD\_NAV\_TAKEOFF}'' became ``the drone is waiting for a message that never comes''); E hit after five drafts because my first three were in my voice, not Demetro's; F hit because the complaint letter had to survive a non-engineering reader. The lesson: communication effectiveness isn't about clarity---it's about \emph{cost asymmetry}. I spent roughly twenty hours on the Markdown doc and Robin spent four integrating; pair walkthroughs inverted that ratio. The same author-hour budget produces dramatically different behaviour change depending on which method the hours go into.

The quieter half of 4b was infrastructure for the team's memory rather than its code. \texttt{brain-dump.md} captured every decision so future sessions didn't start cold. \texttt{MEMORY.md} tracked config values, IP addresses, and baud rates that would otherwise live in one person's head---mine. The blueprint system mapped line ranges for every file over 500 lines so Robin and Demetro could navigate \texttt{vision.py} via a map file instead of asking me to walk them through it. And \texttt{FIELD\_QUICK\_REF.md}, written on the first cancelled flight day, contained every command the team needed to run on the Pi without internet---copy-paste ready, tested. Robin used the blueprint file on at least two occasions during wk21 integration; I know because his commits touched functions inside the line ranges I'd documented, and he didn't message me before either edit. That's a small thing, but it's the proof that my documentation actually reduced my own bottleneck status rather than merely describing it.

\textbf{Self-development programme---what I implemented, not just planned.} Across the project timeline I can identify three implemented development actions, not aspirations. First, after the wk17 FSM conflict I adopted a 24-hour cooling period before responding to technical disagreements---and I held it for every subsequent disagreement, including the complaint letter, which went through six drafts over three days instead of being sent the night I wrote it. Second, after noticing my solo-coding pattern in wk18, I committed to producing at least one handoff artefact per module: the \texttt{--robin} flag, the \texttt{VISION\_PIPELINE} doc, the \texttt{FIELD\_QUICK\_REF}. Against the criterion of whether a teammate could use the module independently within 48 hours, two of three artefacts passed (Robin's integration, pilot's abort usage) and one failed (the 3000-word doc Robin didn't read until forced). Third, I scored my own architecture at 3.0/5.0 on a modularity audit and published the score---not because the number mattered, but because admitting publicly that \texttt{main.py} was my weakest module opened the door for Robin's refactor pushback to land as help rather than criticism.

\textbf{What I leave believing.} I came as a Ukrainian-trained engineer whose reflex under failure was ``one person delivers.'' Graham's Professional Practice session said it plainly: ``it's generally not a lack of resources or technical skill that prevents most teams from succeeding, but is often more related to a set of basic human behaviours.'' I leave believing its near-inversion---\emph{the engineer who delivers alone is one recruitment away from being the single point of failure in whatever she builds}. The wk17 FSM conflict and Robin's wk20 unaided handoff told me, in different registers, that the same work which delivered the report was the work that kept my teammates standing at the edge of their own project. I don't yet know how to hold both truths---that the solo output was necessary \emph{and} that it was damaging---and I suspect that tension won't resolve cleanly. Falsifiable internship signal: in my first six weeks, does at least one teammate modify a module I own without asking permission? If yes, even once, I've started to become an engineer whose work enables other people's work. That is the only kind of engineering I now think is worth aspiring to---and the fact that I arrived at this conviction through failure, not through reading about it, is the thing about this project I will carry longest.
